% Generated by src/43_strobe_checklist.py. Do not edit by hand: every
% anchor here was resolved against manuscript.tex when this was written,
% and editing the table breaks that guarantee silently.
\begingroup
\footnotesize
\setlength{\LTcapwidth}{\textwidth}
\begin{longtable}{@{}p{0.05\textwidth}p{0.30\textwidth}p{0.57\textwidth}@{}}
\caption{STROBE checklist for cohort studies, completed for this manuscript. Each item names the location in the manuscript that addresses it; the locations are resolved against the manuscript source by \texttt{src/43\_strobe\_checklist.py}, which fails rather than emit an item pointing at text that is not there.}\label{tab:supp_strobe}\\
\toprule
Item & Recommendation & Where addressed \\
\midrule
\endfirsthead
\toprule
Item & Recommendation & Where addressed \\
\midrule
\endhead
\bottomrule
\endfoot
1a & Indicate the study's design with a commonly used term in the title or the abstract & The abstract opens the methods sentence with the design term. The title names the populations and the measurement rather than the design, so the abstract carries item 1a. \\
\addlinespace[2pt]
1b & Provide in the abstract an informative and balanced summary of what was done and what was found & The unstructured abstract states aims, design, cohort size, the replication cohort, the primary estimate with its interval, the null truncation result, the audit, and the league-level range including the two leagues where the remedy fails. \\
\addlinespace[2pt]
2 & Explain the scientific background and rationale for the investigation being reported & The Introduction sets out that per-1000-hour rates are the field's convention, that the numerator has been examined and the denominator far less, that prior work on the denominator asked whether exposure time is measured correctly rather than whether it is neutral, and that the offset is a constrained covariate rather than a neutral scaling. \\
\addlinespace[2pt]
3 & State specific objectives including any prespecified hypotheses & The aims sentence appears in the Introduction and in the abstract. No hypothesis was prespecified; the manuscript states this and labels the reference association a worked example. \\
\addlinespace[2pt]
4 & Present key elements of study design early in the paper & Retrospective observational cohort on public appearance records, with a separate appearance-only replication that uses no outcome data. \\
\addlinespace[2pt]
5 & Describe the setting, locations, and relevant dates, including periods of recruitment, exposure, follow-up, and data collection & Fifteen European first-tier competitions; the reference cohort covers English Premier League seasons 2017 to 2025. Provider snapshot dates are recorded in the archived code. \\
\addlinespace[2pt]
6a & Give the eligibility criteria and the sources and methods of selection of participants & Appearances with a recorded date, player identifier and recorded minutes; no recruitment took place, because the study reuses pre-existing public records. \\
\addlinespace[2pt]
7 & Clearly define all outcomes, exposures, predictors, potential confounders, and effect modifiers & Exposure is prior seven-day minutes per 90; the outcome is a same-day event record; the denominator gradient is the slope of log recorded minutes on that exposure. \\
\addlinespace[2pt]
8 & For each variable of interest give sources of data and details of methods of assessment & Two independent providers supply appearances. The injury record is audited rather than assumed, with one sampling frame and one parser applied to both populations. \\
\addlinespace[2pt]
9 & Describe any efforts to address potential sources of bias & The two contaminating defects are separated and measured. Outcome truncation is estimated directly, and independent sources are used to check attribution. Differential ascertainment across squad roles is measured rather than assumed, because the remedy the paper recommends is restriction to one of those roles. \\
\addlinespace[2pt]
10 & Explain how the study size was arrived at & No target size was set. The cohort is every admitted appearance in the selected competitions and seasons, so the counts are a consequence of that frame rather than of a power calculation. The precision actually achieved is reported in place of a power statement, league by league and smallest panel first. \\
\addlinespace[2pt]
11 & Explain how quantitative variables were handled in the analyses, and why groupings were chosen & Recorded minutes enter on the log scale, exposure as minutes per 90. Squad role is a stratifier, not a grouping of a continuous variable. \\
\addlinespace[2pt]
12a & Describe all statistical methods including those used to control for confounding & Ordinary least squares with standard errors clustered on player identifier, plus calendar-phase terms. The gradient is refitted under club-season and fixture clusterings as a sensitivity check, because squad rotation is a club decision and appearances repeat within fixture as well as within player. \\
\addlinespace[2pt]
12b & Describe any methods used to examine subgroups and interactions & Squad role strata are reported separately, with intervals, and the disjointness of the starter interval from the others is stated rather than implied. \\
\addlinespace[2pt]
12c & Explain how missing data were addressed & League-seasons are admitted only above a coverage threshold; the admitted set and its coverage are deposited. No imputation is used in the primary estimate. \\
\addlinespace[2pt]
12d & If applicable, explain how loss to follow-up was addressed & Not applicable: the design has no follow-up of individuals beyond the appearance record itself. \\
\addlinespace[2pt]
12e & Describe any sensitivity analyses & The same league is fitted from two independent providers; a placebo exposure reproduces the whole decomposition; the negative-control analysis is Figure 3; the first-order identity is calibrated against the attenuation observed by sweeping a recorded-minute floor across the range of the gradient; the gradient is refitted under three clusterings; and the full model output for all thirty gradient fits is Appendix A. \\
\addlinespace[2pt]
13 & Report numbers of individuals at each stage of study, and reasons for non-participation & Appearance and player counts are reported for the reference cohort and for each replication population, with the admitted league-seasons deposited. \\
\addlinespace[2pt]
14 & Give characteristics of study participants and information on exposures and potential confounders & Cohort composition and recorded-minute distributions are reported, split by squad role where the split is load-bearing. \\
\addlinespace[2pt]
15 & Report numbers of outcome events or summary measures over time & Event counts are reported for the pooled cohort and for each squad-role stratum. \\
\addlinespace[2pt]
16 & Give unadjusted estimates and, if applicable, confounder-adjusted estimates and their precision & Every estimate is reported with a 95\% interval, and every estimate is named by the model that produced it: Poisson fits are reported as rate ratios and the per-appearance logistic fit as an odds ratio. The primary comparison changes only the offset, so the two forms are reported side by side. \\
\addlinespace[2pt]
17 & Report other analyses done & The cross-league and cross-sex replication and the matched audit of the two injury records are reported, including the two leagues where the recommended remedy does not work. \\
\addlinespace[2pt]
18 & Summarise key results with reference to study objectives & The Discussion opens on the measurement and what it costs a per-minute rate, in the order the aims were stated. \\
\addlinespace[2pt]
19 & Discuss limitations, taking into account sources of potential bias or imprecision & Limitations state the absence of pre-registration, the observational design, the provider dependence, and the audit's non-random sampling frame. \\
\addlinespace[2pt]
20 & Give a cautious overall interpretation considering objectives, limitations, and multiplicity & The interpretation is confined to the estimator's arithmetic. The manuscript states that the gradient is not a biological effect and that no player-level clearance follows from it. \\
\addlinespace[2pt]
21 & Discuss the generalisability of the study results & Generalisability is measured rather than argued: fifteen leagues, two sexes and two providers, with the boundary of the remedy reported where it fails. \\
\addlinespace[2pt]
22 & Give the source of funding and the role of the funders & No specific grant was received from any funding agency in the public, commercial, or not-for-profit sectors. \\
\addlinespace[2pt]
\end{longtable}
\endgroup
